\section{Downstream appendix atlas and claim-type discipline}
\label{app:downstream-atlas-overview}

\begin{fixedbox}
\textbf{Consequence-facing only.} This appendix and every appendix that follows it are consequence-facing only. None functions as a proof source for Sections~1--9. Each is conditional on the fixed closure theorem of the main text, and deleting the entire atlas leaves the foundational arc unchanged.
\end{fixedbox}

\subsection*{Scope and dependency declaration}
This appendix records the governing rule for the Zenodo-only atlas. It depends only on the claim-type discipline of Section~\ref{sec:consequence-pathways}, the AMetric and standing results of Sections~4--5, the operator and physics closure results of Sections~6--7, and the meta-closure results of Sections~8--9. It introduces no new primitives, no new admissibility layers, and no new same-scope selector.

\subsection*{Claim type}
This is an architecture appendix. Its task is not to prove a new theorem family beyond the main text, but to type every later downstream appendix in advance so that no appendix can drift into hidden-proof-source status.

\begin{definition}[Downstream exhaustion schema]
\label{def:downstream-exhaustion-schema}
Fix a same-scope downstream problem posed inside the closure regime of Sections~1--9. A proposed downstream resolution belongs to exactly one of the following classes:
\begin{enumerate}
  \item a boundary selector or admissibility-relevant parameterization;
  \item a same-scope carrier, auxiliary authorizer, or privileged continuation relation;
  \item an operator or mechanism upgrade that treats standing-bearing compatibility as if it were generated later in the interior;
  \item an explicit scope change by added witness load, new ontology, new dynamics, or new target gating; or
  \item a bookkeeping-only redescription that preserves the already fixed same-scope standing-bearing content.
\end{enumerate}
\end{definition}

\begin{theorem}[Downstream exhaustion schema]
\label{thm:downstream-exhaustion-schema}
For any fixed-domain downstream problem posed on the original scope, every same-scope proposed resolution falls under exactly one class of Definition~\ref{def:downstream-exhaustion-schema}. On the fixed scope, only bookkeeping-preserving or quotient-forced consequence survives; the first four classes are either already blocked by Sections~4--9 or else constitute scope change.
\end{theorem}

\begin{proof}
If a proposal acts before the standing-bearing content of the problem is fixed, it is a boundary selector or parameterization and is blocked by the AMetric boundary and the no-second-classifier results of Sections~4--5. If it adds a same-scope authorizer beneath or alongside the unique admissibility gate, it is a carrier or privileged continuation relation and is blocked by the no-carriers, no-repair, and quotient-closure results of Section~5. If it treats quotient-level compatibility as a later mechanism, it violates the operator-exhaustion and role-discipline results of Sections~6--7. If it adds new witness load, new ontology, new dynamics, or new target gating, it changes scope rather than resolving the original problem. If it does none of these, then it preserves the already fixed same-scope standing-bearing content and is bookkeeping-only. No sixth same-scope class remains.
\end{proof}

\begin{corollary}[Atlas typing rule]
Every appendix that follows must therefore do three things explicitly: name its fixed-domain object, state which clause of the downstream exhaustion schema eliminates the prohibited readings, and then state only the bookkeeping-preserving or quotient-forced consequence that survives on the original scope.
\end{corollary}

\begin{proof}
Without a fixed-domain object, the later appendix would leave scope drift uncontrolled. Without an explicit elimination map, the appendix could be misread as freeform interpretation rather than closure consequence. And without a surviving same-scope remainder, the appendix would either become a hidden constructive proposal or a merely rhetorical gloss. The atlas typing rule therefore follows immediately from Theorem~\ref{thm:downstream-exhaustion-schema} and Proposition~\ref{prop:zenodo-appendix-atlas-rule}.
\end{proof}

\begin{reviewbox}
\textbf{No-backflow clause for the atlas.} The atlas records what the fixed closure theorem does to downstream domains. It does not repair, extend, or motivate the theorem chain of Sections~1--9, and it may be removed in its entirety without altering any upstream theorem or proof burden.
\end{reviewbox}
